\section{What the application does}
The Streamlit application, \texttt{app.py}, offers a small interface in which the user chooses a machine type, one of the two models and one of the three scoring methods, uploads a WAV file and presses \emph{Analyse}. The application loads the corresponding checkpoint and Mahalanobis statistics from the handover folder, computes the score with \texttt{dcase\_part1\_lib}, reads the matching threshold from the Part~2 results and displays the score, the threshold, the number of analysed windows and a NORMAL or ABNORMAL decision. Because it imports the same module and the same files as the evaluation scripts, a recording analysed in the application receives exactly the score and decision recorded for it in the prediction files. The interface also shows a short caption warning that the thresholds were calibrated on training recordings and may be optimistic.

\section{Verification status}
During the final audit the application's inference path, from library import through checkpoint and statistics loading to scoring and threshold comparison, was executed programmatically on a WAV file for the Fan and ToyCar models, with both model variants and all three scoring methods. Every configuration produced finite scores and a decision; the first model load took about ten seconds and later analyses well under a second. That test used a synthetically generated noise recording, which every configuration flagged as abnormal, as one would hope. A test of the interactive interface on genuine recordings is not part of the evidence in the submission package; the execution instructions refer to screenshots kept in a separate shared document. The honest description of the application is therefore \emph{implemented and inference-verified}, not user-tested.
